\documentclass[11pt]{article}
\usepackage[T1]{fontenc}
\usepackage[margin=1in]{geometry}
\usepackage{amsmath,amssymb,amsthm}
\usepackage[hidelinks]{hyperref}
\newtheorem{theorem}{Theorem}
\title{An effective lower bound for integer programming with an all-ones right-hand side}
\author{Atlas\thanks{The proof in this note was found using ChatGPT.}}
\date{18 November 2026}
\begin{document}
\maketitle

Consider a binary matrix \(A\) with \(k\) rows and an unrestricted
number of columns. We ask whether \(Ax=\mathbf 1\) for some vector
\(x\) of nonnegative integers, where \(\mathbf 1\) has \(k\) entries.
Write \(N\) for the complete binary input length. We use the incidence
matrix from Lin's reduction to Set Cover~\cite[Lemma~14]{lin}.
Its rows can enforce exact agreement among the edges chosen for the
different pairs of colors.

\hypertarget{ilp-equations-w1-effective}{}
\begin{theorem}\label{ilp-equations-w1-effective}
Assume deterministic \(\mathrm{FPT}\ne\mathrm{W[1]}\). For every fixed
\(C\ge0\) and every finite nonnegative function \(h\) admitting a
computable threshold \(T(j)\) with \(h(k)\le k/j\) for every positive
integer \(j\) and every \(k\ge T(j)\), no deterministic algorithm
decides feasibility of \(Ax=\mathbf 1\) over the nonnegative integers
in \(O(2^{h(k)}N^C)\) time. This holds even when every column of \(A\)
is nonzero, so every feasible coordinate is at most one.
\end{theorem}

\begin{proof}
Reduce from Multicolored Clique parameterized by the number \(\ell\)
of color classes~\cite{cfk}. Decide trivial cases directly, including
an empty class or a pair of classes with no edge between them.
We may then assume \(n\ge\ell\ge2\), where \(n\) is the number of
vertices. Give each vertex a distinct \(s\)-bit name, where
\(s=\lceil\log_2 n\rceil\). For each color \(i\), fix a bijection
\(\sigma_i\) from the other colors to \([\ell-1]\).

For each color \(i\), each function \(f:\{0,1\}\to[\ell-1]\), and each
bit position \(t\in[s]\), make a row \((i,f,t)\). Make one column for
each edge between two color classes. Its entry in row \((i,f,t)\) is
one exactly when the edge has an endpoint \(v\) of color \(i\), its
other endpoint has color \(j\), and \(f(\operatorname{bit}_t(v))=
\sigma_i(j)\), where \(\operatorname{bit}_t(v)\) is bit \(t\) of
the name of \(v\). This is Lin's incidence construction, with the
bit taken at the endpoint of the row's color. Every column has a
one in a row for a constant function at either endpoint's color.

A multicolored clique gives a solution by setting its edge variables
to one and all others to zero. In a row \((i,f,t)\), every selected
edge incident with color \(i\) shares the clique vertex of that color.
The function value on its bit picks exactly one partner color. Thus
exactly one selected edge contributes to the row.

Conversely, fix two colors \(i\) and \(j\). A row for color \(i\) and
the function constantly equal to \(\sigma_i(j)\) sums precisely the
variables for edges between these two colors. Since the sum is one
and the coordinates are nonnegative integers, exactly one such
edge is selected. This proves that every coordinate is binary. If
two selected edges from \(i\) to distinct colors \(j\) and \(j'\)
use different endpoints of color \(i\), choose a bit position where
those endpoint names differ. Map the first bit to \(\sigma_i(j)\)
and the second to \(\sigma_i(j')\). Both edges then contribute to
the resulting row, contradicting its sum of one. Hence all selected
edges at each color have a common endpoint. These endpoints form
a multicolored clique.

There are \(k=\ell(\ell-1)^2s\le 2\ell^3\log_2 n\) rows and at most
\(n^2\) columns. Enumerating the rows, columns, and entries, including
the right-hand side, takes \(F(\ell)n^{O(1)}\) time for a computable
\(F\), which also bounds the complete output length. Suppose the solver
exists and set \(j=2\ell^3\). If \(k\ge T(j)\), then \(h(k)\le\log_2 n\).
Otherwise, fix one integer \(H\ge0\) bounding \(h(i)\) for \(i<T(1)\).
The finite bound \(H\) can be built into the resulting algorithm. It does
not require a procedure for computing \(h\).
For \(i\ge T(1)\), we have \(h(i)\le i\), so in the second case
\(h(k)\le H+T(j)\). In both cases the solver and construction give
deterministic fixed-parameter time in the number of colors, a contradiction.
\end{proof}

Covering inequalities with one cardinality constraint already appear
in the Set Cover formulation of Dvo\v{r}\'ak et al.~\cite[Theorem~25]{degko}.
Rohwedder and W\k{e}grzycki~\cite[Lemmas~15--16]{rw} also use exact
cardinality and complemented incidence in related reductions.
Here the constant-function rows enforce the exact edge choices
directly on Lin's matrix. The all-ones restriction admits a
\(2^kN^{O(1)}\) subset dynamic program. Delete zero columns and
start with the empty row set and repeatedly add a disjoint column
support to an attainable set. An attainable full row set is equivalent to feasibility.
For general right-hand sides, Knop, Pilipczuk, and
Wrochna~\cite[Theorem~1]{kpwconf} give an ETH lower bound at the
\(k\log k\) scale, with column count and maximum right-hand side
\(O(k\log k)\).

\begin{thebibliography}{9}
\bibitem{cfk}
Marek Cygan et al.
\newblock \emph{Parameterized Algorithms}.
\newblock Springer, 2015.
\newblock \href{https://doi.org/10.1007/978-3-319-21275-3}{doi:10.1007/978-3-319-21275-3}.

\bibitem{lin}
Bingkai Lin.
\newblock A simple gap-producing reduction for the parameterized Set Cover problem.
\newblock In \emph{ICALP}, LIPIcs 132, 81:1--81:15, 2019.
\newblock \href{https://doi.org/10.4230/LIPIcs.ICALP.2019.81}{doi:10.4230/LIPIcs.ICALP.2019.81}.

\bibitem{kpwconf}
Du\v{s}an Knop, Micha{\l} Pilipczuk, and Marcin Wrochna.
\newblock Tight complexity lower bounds for integer linear programming with few constraints.
\newblock \emph{ACM Transactions on Computation Theory} 12(3), Article 19,
2020, \href{https://doi.org/10.1145/3397484}{doi:10.1145/3397484}.
\newblock The inspected edition is \emph{STACS}, LIPIcs 126, 44:1--44:15,
2019, \href{https://doi.org/10.4230/LIPIcs.STACS.2019.44}{doi:10.4230/LIPIcs.STACS.2019.44}.
\bibitem{degko}
P. Dvo\v{r}\'ak, E. Eiben, R. Ganian, D. Knop, and S. Ordyniak.
The complexity landscape of decompositional parameters for ILP:
Programs with few global variables and constraints.
\emph{Artificial Intelligence} 300, Article 103561, 2021.
\href{https://doi.org/10.1016/j.artint.2021.103561}{DOI}.

\bibitem{rw}
L. Rohwedder and K. W\k{e}grzycki.
Fine-Grained Equivalence for Problems Related to Integer Linear Programming.
\href{https://doi.org/10.4230/LIPIcs.ITCS.2025.83}{ITCS 2025, Article 83}.

\end{thebibliography}
\end{document}
